\section{Exemplos resolvidos}

\begin{frame}{Exemplo 1 — circuito de TUG em área seca}
\small
\begin{columns}[T,onlytextwidth]
\column{0.49\textwidth}
\begin{caixaEx}
Circuito em 220 V, $P=\SI{1100}{VA}$, comprimento de estudo $L=\SI{30}{m}$:
\[
I_b=\frac{1100}{220}=\SI{5.00}{A}.
\]
Adota-se inicialmente Cu 2,5 mm$^2$ e disjuntor 16 A.
\end{caixaEx}
\column{0.49\textwidth}
\textbf{A seleção só fecha após verificar:}
\begin{enumerate}
\item método de instalação;
\item temperatura e agrupamento;
\item $I_b\le I_n\le I_z$;
\item queda de tensão;
\item curto e desligamento;
\item necessidade de proteção diferencial.
\end{enumerate}
\end{columns}
\end{frame}

\begin{frame}{Exemplo 2 — circuito de chuveiro}
\small
\begin{columns}[T,onlytextwidth]
\column{0.49\textwidth}
\begin{caixaEx}
$P=\SI{5.5}{kW}$, $V=\SI{220}{V}$:
\[
I_b=\frac{5500}{220}=\SI{25}{A}.
\]
Para este exemplo de 5,5 kW, adota-se inicialmente circuito dedicado, Cu 6 mm$^2$ e DJ 32 A, após as verificações de capacidade de condução e queda. O projeto residencial adiante usa chuveiros de 4,5 kW e outra seleção.
\end{caixaEx}
\column{0.49\textwidth}
\begin{caixaNorma}[title=Proteção adicional]
O circuito que serve ponto de utilização em local contendo chuveiro deve receber proteção diferencial-residual de alta sensibilidade conforme os requisitos aplicáveis.
\end{caixaNorma}
\begin{caixaAt}
Não escolher seção apenas por uma tabela de potência de chuveiro. Comprimento, método, temperatura e agrupamento mudam o resultado.
\end{caixaAt}
\end{columns}
\end{frame}

\begin{frame}{Exemplo 3 — balanceamento antes do padrão de entrada}
\small
\begin{columns}[T,onlytextwidth]
\column{0.52\textwidth}
No projeto residencial desenvolvido neste material:
\[
S_A=\SI{11.36}{kVA},\qquad S_B=S_C=\SI{11.30}{kVA}.
\]
Em 220 V fase-neutro:
\[
I_A\approx\SI{51.6}{A},\qquad I_B=I_C\approx\SI{51.4}{A}.
\]
\column{0.46\textwidth}
\begin{caixaProjeto}[title=Distribuição]
\begin{itemize}
\item A: C1+C3+C7+C9+C13+C16;
\item B: C2+C4+C10+C12+C15;
\item C: C5+C6+C8+C11+C14.
\end{itemize}
A soma fecha com os 16 circuitos e mantém um chuveiro em cada fase. Este é o balanceamento dos valores instalados; cenários de demanda também devem ser verificados.
\end{caixaProjeto}
\end{columns}
\end{frame}

\begin{frame}{Exemplo 4 — enquadramento Equatorial Maranhão}
\small
\begin{columns}[T,onlytextwidth]
\column{0.50\textwidth}
Carga instalada do projeto:
\[
P_{inst,eq}=\SI{33.96}{kW}.
\]
A NT.00001.EQTL Rev.09 enquadra a faixa de 24,1 a 38 kW, na Tabela 1 de 220/380 V, em ligação trifásica com disjuntor termomagnético de entrada de 63 A.
\column{0.48\textwidth}
\begin{caixaNorma}[title=No Maranhão]
A ligação trifásica urbana é 380/220 V, 3F+N. A mesma tabela indica 10 mm$^2$ de cobre como mínimo do condutor do cliente no padrão correspondente.
\end{caixaNorma}
\begin{caixaAt}
O alimentador interno é recalculado pela NBR 5410; não se copia automaticamente o mínimo do padrão da concessionária.
\end{caixaAt}
\end{columns}
\end{frame}

\begin{frame}{Exemplo 5 — recarga de VE e aumento de carga}
\small
\begin{columns}[T,onlytextwidth]
\column{0.49\textwidth}
Se ao projeto de $\SI{33.96}{kW}$ for acrescentado um SAVE de $\SI{7.4}{kW}$:
\[
P_{inst,novo}=33{,}96+7{,}40=\SI{41.36}{kW}.
\]
O exemplo passa à faixa de 38,1 a 48 kW, associada a 80 A na Tabela 1; o enquadramento da entrada precisa ser revisto.
\column{0.49\textwidth}
\begin{caixaProjeto}[title=Consequência]
Além do novo circuito e da capacidade do QD/alimentador, a conexão da estação de recarga deve observar a NT.00042.EQTL Rev.04 e eventual necessidade de aumento de carga perante a distribuidora.
\end{caixaProjeto}
\end{columns}
\end{frame}

%=====================================================================
